% ============================================================
% 第1章 问题背景与重述
% ============================================================

\section{问题重述}

\subsection{问题的背景}
板凳龙是浙闽地区传统民俗活动，由数十至上百条板凳首尾相连，形成蜿蜒长龙。舞龙时，龙头率先沿螺线盘入，龙身和龙尾依次跟随，整体呈圆盘状。表演中，在能自如盘入和盘出的前提下，所需场地面积越小、行进速度越快，观赏性越高。本题针对由223节板凳组成的舞龙队，龙头板长341 cm，龙身和龙尾板长220 cm，板宽30 cm，各节设有两个孔，孔径5.5 cm，孔中心距板头27.5 cm，相邻板凳通过把手连接。需建立数学模型，研究沿等距螺线运动时各把手的位置与速度、碰撞终止时刻、调头空间及最小螺距、S形调头路径优化以及最大行进速度等问题，为舞龙表演的规划与控制提供理论依据。

\subsection{问题的提出}
某板凳龙由223节板凳组成，其中第1节为龙头，其后221节为龙身，最后1节为龙尾。龙头板长为341 cm，龙身和龙尾板长均为220 cm，板宽均为30 cm。每节板凳上开有两个圆孔，孔径5.5 cm，孔中心距最近的板头27.5 cm，相邻板凳通过把手穿过圆孔连接。舞龙队沿螺距为55 cm的等距螺线顺时针盘入，各把手中心均位于螺线上，龙头前把手行进速度恒为1 m/s，初始时龙头位于螺线第16圈A点。

\textbf{问题1：}建立舞龙队沿给定螺线运动的数学模型，给出从初始时刻到300 s内每秒整个舞龙队各把手中心（包括龙头、龙身和龙尾各前把手及龙尾后把手）的位置和速度，将结果保存至文件result1.xlsx中；并在论文中列出0 s、60 s、120 s、180 s、240 s、300 s时龙头前把手、龙头后第1、51、101、151、201节龙身前把手和龙尾后把手的位置和速度。

\textbf{问题2：}在问题1设定的盘入螺线基础上，确定因相邻板凳发生碰撞而不能再继续盘入的终止时刻，并给出此时上述指定把手的位置和速度，将结果保存至result2.xlsx。

\textbf{问题3：}从盘入到盘出需要调头，调头空间为以螺线中心为圆心、直径为9 m的圆形区域。确定最小螺距，使得龙头前把手能够沿相应的螺线盘入至该调头空间边界，即恰好与调头区域相切。

\textbf{问题4：}盘入螺距为1.7 m，盘出螺线与盘入螺线关于螺线中心呈中心对称，调头路径为由两段圆弧相切连接而成的S形曲线（前一段圆弧半径是后一段的2倍），且与盘入、盘出螺线均相切。研究能否调整两段圆弧（仍保持相切）使调头曲线更短。以调头开始时刻为零时刻，给出从-100 s到100 s内每秒舞龙队的位置和速度，将结果保存至result4.xlsx，并在论文中列出-100 s、-50 s、0.5 s、50 s、100 s时指定把手的位置和速度。

\textbf{问题5：}在问题4设定的完整路径（盘入—调头—盘出）下，龙头行进速度保持不变，确定龙头的最大行进速度，使得舞龙队各把手的速度均不超过2 m/s。

